\documentclass[a4paper,12pt]{article}

\usepackage[utf8]{inputenc}
\usepackage[T5]{fontenc}
\usepackage[vietnamese]{babel}
\usepackage{amsmath, amssymb}
\usepackage{geometry}
\usepackage{enumitem}
\usepackage{tcolorbox}
\usepackage{listings}

\geometry{margin=2cm}

\begin{document}

\begin{center}
    \textbf{\Large BÀI TẬP DFS}
\end{center}

\section*{TỔNG QUAN}
\begin{tcolorbox}
Bài toán sử dụng DFS để duyệt đồ thị và xác định thứ tự thăm các đỉnh.
\end{tcolorbox}

\begin{itemize}
\item Không dùng thuật toán nâng cao
\item Chỉ cần DFS cơ bản
\item Rèn luyện kỹ năng duyệt đồ thị
\end{itemize}

\section*{ĐỀ BÀI}
\begin{tcolorbox}
Cho đồ thị vô hướng gồm $n$ đỉnh và $m$ cạnh.

Hãy thực hiện duyệt DFS bắt đầu từ đỉnh $1$ và in ra thứ tự các đỉnh được thăm.
\end{tcolorbox}

\begin{itemize}
\item Các đỉnh được đánh số từ $1$ đến $n$
\item Nếu có nhiều lựa chọn, ưu tiên đỉnh có số nhỏ hơn
\end{itemize}

\section*{DỮ LIỆU VÀO}
\begin{tcolorbox}
Dòng đầu: $n, m$

$m$ dòng tiếp theo: mỗi dòng gồm $u, v$
\end{tcolorbox}

\begin{itemize}
\item $1 \le n \le 1000$
\item $0 \le m \le \dfrac{n(n-1)}{2}$
\end{itemize}

\section*{DỮ LIỆU RA}
\begin{tcolorbox}
In ra thứ tự DFS bắt đầu từ đỉnh $1$
\end{tcolorbox}

\begin{itemize}
\item Các đỉnh cách nhau bởi dấu cách
\end{itemize}

\section*{VÍ DỤ}

\begin{tcolorbox}[title=Input]
\begin{lstlisting}
5 5
1 2
2 3
1 4
2 5
4 5
\end{lstlisting}
\end{tcolorbox}

\begin{tcolorbox}[title=Output]
\begin{lstlisting}
1 2 3 5 4
\end{lstlisting}
\end{tcolorbox}

\section*{GIẢI THÍCH}
\begin{itemize}
\item Bắt đầu từ 1 → đi 2
\item từ 2 → đi 3
\item quay lại → đi 5
\item từ 5 → đi 4
\end{itemize}

\begin{tcolorbox}
DFS luôn đi sâu trước rồi mới quay lui
\end{tcolorbox}

\section*{GỢI Ý}
\begin{itemize}
\item Dùng vector<int> để lưu danh sách kề
\item Sắp xếp danh sách kề
\item Dùng DFS đệ quy
\item Mảng visited để tránh lặp
\end{itemize}

\begin{tcolorbox}
Nếu không sort, thứ tự DFS sẽ khác
\end{tcolorbox}

\section*{SUBTASK}
\begin{itemize}
\item Subtask 1: $n \le 20$ (30\%)
\item Subtask 2: $n \le 200$ (30\%)
\item Subtask 3: $n \le 1000$ (40\%)
\end{itemize}

\section*{LƯU Ý}
\begin{itemize}
\item DFS phụ thuộc thứ tự duyệt
\item luôn cần visited
\item nên sort adjacency list
\end{itemize}

\begin{tcolorbox}
Đây là bài nền tảng cực quan trọng
\end{tcolorbox}

\end{document}